% !TeX root = ../main.tex

\section{Empirical Study}

This section presents a comprehensive study of the performace of LLMs on real-world Java projects. Without employing complex dependency analysis and other enhancement methods, there are two naive translation schemes:

\textbf{File-level Translation}: File-by-file translation using LLMs with minimal context.

\textbf{Method-level Translation}: Method-by-method translation using LLMs with minimal context.

\subsection{Research Questions}

Our study addresses the following research questions:

\begin{itemize}
\item \textbf{RQ1}: What are the main problems that arise during the translation process using these two methods?
\item \textbf{RQ2}: What is the problems distribution of these two methods?
\item \textbf{RQ3}: How does the performance compare between File-level and Method-level Translation?
\end{itemize}

\subsection{Setup}
\subsubsection{Dataset}
We evaluated the performance of LLMs on file groups from real-world Java projects representing different domains and complexity levels:
\begin{table}[h]
\centering
\small
\begin{tabular}{lcccc}
\toprule
\textbf{Translation Module}\textbf{Main Class} & \textbf{File Count} & \textbf{Method Count} & \textbf{LOC} & \textbf{Domain} \\
\midrule
Module1 & Cookie & 8 & 45 & 1,247 & HTTP Library \\
Module2 & JSONParser & 6 & 32 & 856 & Data Parsing \\
Module3 & ReadRowHolder & 9 & 67 & 2,134 & Excel Processing \\
Module4 & ReadSheetHolder & 13 & 89 & 3,021 & Excel Processing \\
\midrule
\textbf{Total} & 36 & 233 & 7,258 & -- \\
\bottomrule
\end{tabular}
\caption{Characteristics of evaluation dataset}
\label{tab:dataset}
\end{table}

After translation, we manually analyzed the various problems in the results and summarized and statistically analyzed them.
\subsection{Results}

